% Generated draft: requires mathematical and source review.
\documentclass[11pt]{article}
\usepackage[margin=25mm]{geometry}
\usepackage{fontspec}
\setmainfont{Latin Modern Roman}
\usepackage{amsmath,amssymb,mathtools,unicode-math}
\setmathfont{Latin Modern Math}
\usepackage{enumitem,needspace,xurl,hyperref}
\setlength{\parindent}{0pt}
\setlength{\parskip}{5pt}
\setlength{\emergencystretch}{3em}
\newcommand{\N}{\mathbb{N}}
\newcommand{\Z}{\mathbb{Z}}
\newcommand{\Q}{\mathbb{Q}}
\newcommand{\R}{\mathbb{R}}
\newcommand{\C}{\mathbb{C}}
\newcommand{\F}{\mathbb{F}}
\newcommand{\A}{\mathbb{A}}
\newcommand{\PP}{\mathbb{P}}
\newcommand{\E}{\mathbb{E}}
\newcommand{\eps}{\varepsilon}
\newcommand{\dd}{\,\mathrm{d}}
\begin{document}
\section*{UnsolvedMath ID: 20003321}
\textbf{Source problem code:} catalogue entry
\subsection{Definitions and mathematical statement}
The aid recovers Question 35 as an effective Oppenheim problem: let $Q$ be a nondegenerate indefinite quadratic form on $\R^d$, $d\geq3$, irrational in the sense of not being a real scalar multiple of a rational form. For every $\varepsilon>0$, seek an effective bound $T=T_Q(\varepsilon)$ and $x\in\Z^d$ with
\[
 0<|Q(x)|<\varepsilon,\qquad\|x\|<T.
\]
The hypotheses come from preceding source context as reported by the aid, not the locally extracted question alone. The coefficient/Diophantine dependence of $T$ is part of the request; writing $T(\varepsilon)$ does not make it uniform in $Q$. The source's purported easier version omits $0<|Q(x)|$ and any $x\ne0$ condition, so it is literally satisfied by $x=0$; the intended nontrivial version requires a repair. Its reference to upcoming Margulis--G\"otze work for $d\geq5$ is historical.

The aid reports modern effective theorems answering the broad strict question, but gives no full references, explicit general bound or coefficient hypotheses. Its explicit special case is the eight-variable form, for irrational $\alpha>0$,
\[
 Q_\alpha(u,v)=\sum_{i=1}^4u_i^2-\alpha\sum_{i=1}^4v_i^2,
 \qquad u,v\in\Z^4.
\]
Define
\[
 N_\alpha(\varepsilon)=\min\{n\geq1:\|n\alpha\|<\varepsilon\},
 \qquad
 H_\alpha(\varepsilon)=\min\{\|(u,v)\|_\infty:
                  0<|Q_\alpha(u,v)|<\varepsilon\},
\]
where $\|n\alpha\|$ is nearest-integer distance. For $0<\varepsilon<\min\{\alpha,1/2\}$, the reported four-square transfer theorem is
\[
 \frac12\sqrt{N_\alpha(\varepsilon)}
 \leq H_\alpha(\varepsilon)
 \leq\sqrt{(\alpha+1)N_\alpha(\varepsilon)}.
\]
Irrationality ensures an approximation from positive integer square sums is nonzero, and the small-$\varepsilon$ conditions exclude a spurious zero block. For badly approximable $\alpha$, meaning $\inf_{n\geq1}n\|n\alpha\|>0$, the aid obtains
\[
 H_\alpha(\varepsilon)=\Theta_\alpha(\varepsilon^{-1/2})
 \quad(\varepsilon\to0).
\]
This is a coefficient-dependent asymptotic for the stated family, not a sharp bound for all forms. A separate reported irrational perturbation of $x^2+y^2-3z^2$ shows why coefficient-free uniform bounds fail; the perturbation formula is not supplied.

The extraction then begins separate Section 5.6. For a nondegenerate indefinite rational binary form $Q(x_1,x_2)=ax_1^2+bx_1x_2+cx_2^2$ that represents no nonzero rational zero, it defines
\[
 m(Q,x)=\inf_{z\in\Z^2}|Q(x+z)|,\qquad
 m(Q)=\sup_{x\in\R^2}m(Q,x),
\]
and, if the supremum is isolated,
\[
 m_2(Q)=\sup\{m(Q,x):m(Q,x)<m(Q)\}.
\]
It motivates $m(Q)$ via norm forms of real quadratic fields in an integral basis and states an equivalence with a Euclidean algorithm when $m(Q)<1$. The precise intended assertion concerns the norm-Euclidean condition, not every possible Euclidean function; normalization and any attainment/boundary qualifications need its actual reference. It then introduces a conjecture communicated by A. Pollington but does not give it. This binary inhomogeneous material and the page header are not clauses of the effective higher-dimensional Question 35.
\subsection{Short English statement}
The strict small-value question needs coefficient-dependent effective bounds and excludes the zero vector automatically. The aid gives a precise eight-variable height estimate and badly-approximable asymptotic, while general effective references are missing. The binary inhomogeneous minima and norm-Euclidean discussion belong to the next section.
\subsection{Sources}
\begin{itemize}
\item Local UnsolvedMath registry, permanent ID 20003321, source code catalogue entry; \texttt{aim-other-notes.json}, notes[127], source tag question.
\item AIM, \emph{Open Problems from the Workshop Emerging Applications of Measure Rigidity}, 2004, Section 5.5, Question 35, G. Margulis [130], followed by Section 5.6: \url{https://aimath.org/WWN/measrigid/measrigid.pdf}. No [130] bibliographic entry is supplied.
\item G. Margulis and F. G\"otze's then-upcoming $d\geq5$ work is named in the source. Modern effective Oppenheim results are claimed by the aid without full citations. The following missing conjecture is attributed to A. Pollington.
\item Machine-generated research attempt 1: partial result, solved status at the run, rigor labelled ``proved,'' low novelty confidence; registry solved/``SOLVED-IN-LITERATURE.'' Not peer reviewed; independently verify before citation. Registry updated 14 August 2026.
\end{itemize}
\end{document}
